\section{Calibration}
\label{sec:calibration}

We calibrate the model to empirically grounded parameter values from the mutual fund and portfolio choice literature, ensuring that the implied equilibrium outcomes fall within empirically plausible ranges for delegated portfolio management.

We set \(\gamma=5\) as the baseline investor risk aversion, with sensitivity analysis spanning \(\gamma\in[2,10]\). These values are consistent with standard portfolio-choice calibrations \citep[][]{campbell2002strategic} and reflect the observed heterogeneity in risk preferences across investors \citep[][]{guiso2018time}.

We set \(\eta=15\) as the baseline manager risk aversion over fee revenue, with sensitivity spanning \(\eta\in[3,30]\). The key economic rationale is that fee income is the manager's primary and largely undiversifiable source of compensation. A portfolio manager whose livelihood depends on AUM-linked revenue cannot hedge fee-income risk across multiple positions the way a diversified investor can. Career concerns further amplify effective risk aversion: a manager who underperforms relative to peers faces redemptions, reputational damage, and possible termination, all of which are convex in downside outcomes \citep[][]{chevalier1997risk, BrownHarlowStarks1996}. Setting \(\eta=15\) --- three times the investor's risk aversion \(\gamma=5\) --- reflects this concentration premium over fee-revenue risk. The sensitivity analysis in Section~\ref{sec:numerical} confirms that the model's qualitative predictions are robust across the full sweep of \(\eta\in[3,30]\).

We set \(\Sharpe=0.35\) (annualized) for an active mutual fund baseline, consistent with the performance range documented for skilled managers net of transaction costs \citep[][]{carhart1997persistence}. We explore \(\Sharpe\in[0.10,1.00]\) in sensitivity analysis, with higher values (e.g., \(\Sharpe=0.60\)) representing particularly skilled managers or hedge fund strategies \citep[][]{fung2004hedge}.

We use \(r_f=3.70\%\) (annual), based on the 3-month U.S. Treasury constant maturity rate as of January 2026.\footnote{Source: FRED series DGS3MO, 2026-01-23: 3.70\%. For euro-denominated settings, the ECB €STR (euro short-term rate) was approximately 1.93\% as of the same date.} The qualitative results are robust to alternative risk-free rate values across plausible ranges, as $r_f$ enters the model only through the flow threshold and the net return level, and its effects on equilibrium $\sigma_d^*$ and $\phi^*$ are second-order relative to the convexity and risk-aversion parameters.

Empirical studies document strong convexity in the flow-performance relationship, whereby top-performing funds attract disproportionately large inflows \citep[][]{sirri1998costly, ferreira2012flow}. To match realistic annual flow responses, we set $f_1=1.5$ and $f_2=25$, where excess returns are measured in decimal units (e.g., a 5\% excess return corresponds to 0.05). This parameterization implies that a fund with \(+5\%\) excess return experiences flows of approximately \(+13.8\%\) of AUM, while a fund with \(+10\%\) excess return sees flows around \(+40\%\), capturing the strong convexity observed in empirical flow--performance relationships. We explore \(f_1\in[0.5,3]\) and \(f_2\in[5,60]\) to assess sensitivity to flow convexity.

We normalize \(A_0=1\) since the problem is scale-invariant with proportional fees and homogeneous preferences. If one were to incorporate decreasing returns to scale or capacity constraints, AUM scale would interact with fund performance \citep[][]{berk2004mutual}, but we abstract from such effects in our baseline specification.

With these parameters, the model produces optimal volatility targets \(\sigma_d^*\in[8\%,20\%]\) (annualized), which are consistent with observed volatility levels for actively managed equity mutual funds. For reference, \citet[][]{barras2010false} use a monthly market-return standard deviation of approximately 4.6\%, corresponding to roughly 16\% annualized. Individual fund volatilities can be higher or lower depending on style and leverage, but our target range aligns with typical volatilities of delegated equity portfolios.
